%------------------------------------------------------------------------------%
\section{Convergência Absoluta e Condicional}

Testar diretamente a convergência de séries com termos positivos e
negativos arbitrários precisa considerar o arranjo particular de sinais,
assim não temos testes gerais para esses casos.
A solução que empregaremos é comparar as séries com termos quaisquer
com séries com termos não negativos.
A definição a seguir define os conceitos de convergência condicional e absoluta
que empregaremos para esse fim.

%------------------------------------------------------------------------------%
\begin{definicao}{Convergência Absoluta e Condicional}{def:convergencia-absoluta}
  Dizemos que a série
  \(
    \sum_{n=1}^\infty a_n
  \)
  \emph{Converge Absolutamente}%
  \index{Série!convergência absoluta}%
  \index{Série!convergência condicional},
  ou é \emph{Absolutamente Convergente}, se a série
  com os valores absolutos
  \(
    \sum_{n=1}^\infty \,\abs{a_n}
  \)
  converge.

  Uma série que converge mas não é absolutamente convergente
  é chamada \emph{Condicionalmente Convergente}.
\end{definicao}
%------------------------------------------------------------------------------%

Quando queremos verificar a convergência de uma série com termos positivos
e negativos podemos construir a série que soma os módulos dos termos dessa
série e verificar a convergência dessa nova série usando um dos testes
para as séries de termos não negativos.
Se essa série convergir usamos o teorema a seguir para garantir
que a série original também converge.

%------------------------------------------------------------------------------%
\begin{teorema}{Teste da Convergência Absoluta}{teo:teste-convergencia-absoluta}
  Uma\index{Teste!convergência absoluta}
  série absolutamente convergente é convergente.
\end{teorema}
%------------------------------------------------------------------------------%
\begin{proof}
  Seja
  \[
    \sum_{n=1}^\infty a_n
  \]
  uma série absolutamente convergente. Observando que
  \[
    -\abs{a_n} \leq a_n \leq \abs{a_n}
  \]
  e somando $\abs{a_n}$ em todos os termos temos
  \[
    0 \leq a_n + \abs{a_n} \leq 2\abs{a_n}    \,.
  \]
  Como a série é absolutamente convergente, segue que
  \[
    \sum_{n=1}^\infty 2\abs{a_n}
  \]
  converge.
  Assim, pelo Teste da Comparação a série
  \[
    \sum_{n=1}^\infty \big( a_n + \abs{a_n} \big)
  \]
  converge.
  Consequentemente, a série
  \[
    \sum_{n=1}^\infty a_n
  \]
  converge, pois
  \[
    \sum_{n=1}^\infty a_n = \sum_{n=1}^\infty \big( a_n + \abs{a_n} \big) - \abs{a_n}
  \]
  e as séries
  \[
    \sum_{n=1}^\infty \big( a_n + \abs{a_n} \big)
    \qquad\text{ e }\qquad
    \sum_{n=1}^\infty \abs{a_n}
  \]
  convergem.
\end{proof}
%------------------------------------------------------------------------------%

Note que não é possível afirmar que uma série que converge também converge
absolutamente. Como exemplo, observe a série~$p$ alternada
\[
  \sum_{n=1}^\infty \frac{\,(-1)^{n+1}\,}{n^p}       \,.
\]
Essa série converge se $0<p\leq 1$,
entretanto esta série não é absolutamente convergente, pois
\[
  \sum_{n=1}^\infty \abs{\frac{(-1)^{n+1}}{n^p}} = \sum_{n=1}^\infty \frac{1}{n^p}
\]
e a série~$p$ não converge se $p\leq 1$,
ou seja, a série~$p$ alternada, com $0<p\leq 1$, é condicionalmente convergente.

Observe que dada uma série convergente de termos não negativos $a_n$, então,
a série alternada
\[
  \sum (-1)^{n+1}a_n
\]
é absolutamente convergente.
Note que, todas as séries convergentes com termo geral não negativo que vimos
anteriormente, tem sua respectiva série alternada absolutamente convergente.

Os exemplos a seguir mostram como usar o teste da convergência absoluta.

%------------------------------------------------------------------------------%
\begin{example}
  Use o teste da convergência absoluta para mostrar que
  \[
    \sum_{n=1}^\infty (-1)^n\frac{(n!)^2\,3^n}{(2n+1)!}
  \]
  é uma série convergente.
  \solution
  Mostraremos que a série é absolutamente convergente, ou seja,
  que a série
  \[
    \sum_{n=1}^\infty\abs{a_n} = \sum_{n=1}^\infty \frac{(n!)^23^n}{(2n+1)!}
  \]
  converge.
  Para isso aplicaremos o Teste da Razão
  \begin{align*}
  	\frac{\abs{a_{n+1}}}{\abs{a_n}}
    & = \frac{\big((n+1)!\big)^2\, 3^{n+1}}{(2n+3)!} \; \frac{(2n+1)!}{(n!)^2\,3^n}     \\[3mm]
  	& = \frac{\,(n+1)^2 (n!)^2\; 3^{n} 3 \; (2n+1)!\,}
             {\,(2n+3)(2n+2)(2n+1)!\; (n!)^2\; 3^n\,} \\[3mm]
  	& = \frac{3(n+1)^{2}}{\,2{(n+1)}(2n+3)\,}         \\[3mm]
  	& = \frac{3(n+1)}{\,2(2n+3)\,}
  \end{align*}
  Assim, aplicando L'Hôpital, segue que
  \[
    \lim \frac{\,\abs{a_{n+1}}\,}{\abs{a_n}}
      = \lim \frac{3(n+1)}{\,2(2n+3)\,}
      = \frac{3}{\,4\,}
      < 1   \;.
  \]
  Portanto, pelo Teste da Razão a série $\sum\abs{a_n}$ é convergente.
  Então, pelo Teste da Convergência Absoluta, concluímos que a série
  \[
    \sum_{n=1}^\infty (-1)^n\frac{(n!)^23^n}{(2n+1)!}
  \]
  é absolutamente convergente.
\end{example}
%------------------------------------------------------------------------------%

%------------------------------------------------------------------------------%
\begin{example}
  Verifique se a série
  \(\;
    \sum_{n=1}^\infty (-1)^n \frac{\,\arctan(n)\,}{n^2+1}
  \;\)
  é absolutamente convergente.
  \solution
  Temos que verificar a convergência da série
  \[
    \sum_{n=1}^\infty \frac{\,\arctan(n)\,}{n^2+1}  \;.
  \]
  Pelo Teste da Integral sabemos que essa série converge se, e somente se,
  a integral
  \[
    \int_1^\infty\frac{\,\arctan(x)\,}{x^2+1}\,dx
  \]
  for convergente.
  Para calcular a integral usamos a mudança de variáveis
  \[
    u = \arctan(x)
    \qquad\qquad
    du = \frac{dx}{x^2+1}
  \]
  que leva aos novos limites de integração
  \[
    \arctan(1)=\frac{\pi}{4}
    \qquad\qquad
    \lim_{x\to \infty} \arctan(x) = \frac{\pi}{2}
  \]
  segue que
  \begin{align*}
    \int_1^\infty \frac{\,\arctan(x)\,}{x^2+1}\,dx
    & = \int_{\sfrac{\pi}{4}}^{\sfrac{\pi}{2}} u\,du             \\[3mm]
    & = \frac{u^2}{2}\evalat{\sfrac{\pi}{4}}{\sfrac{\pi}{2}}     \\[3mm]
    & = \frac{1}{2}\left(\frac{\pi^2}{4}-\frac{\pi^2}{16}\right) \\[3mm]
    & = \frac{3\pi^2}{32}  \;.
  \end{align*}
  Como a integral é convergente a série é absolutamente convergente.
\end{example}
%------------------------------------------------------------------------------%

%------------------------------------------------------------------------------%
